\documentclass{beamer}
\usepackage[spanish]{babel}
\usepackage{graphicx}
\usepackage{booktabs}

\usetheme{Madrid}

\title{Fase 3: Entrenamiento y Modelado Predictivo}
\subtitle{Sistema AASIM}
\author{Ingeniería TI}
\date{\today}

\begin{document}

\begin{frame}
    \titlepage
\end{frame}

\begin{frame}{Objetivos de la Fase 3}
    \begin{itemize}
        \item Consolidar la Tabla Base Analítica (ABT).
        \item Establecer pipelines formales y control de particiones anti-fuga.
        \item Comparar Baselines (Ridge, SVR) frente a Avanzados (LightGBM, PyTorch).
        \item Aislar y evaluar el poder predictivo (Ablation).
    \end{itemize}
\end{frame}

\begin{frame}{Resultados de los Baselines}
    \begin{block}{Observación Crítica}
    Con una muestra de validación pequeña ($n_{train} \approx 691$ ventanas), el modelo \textbf{SVR} superó la prueba logrando el mejor desempeño.
    \end{block}
    
    \begin{itemize}
        \item \textbf{SVR:} MAE = $14.69$ mmHg
        \item \textbf{Dummy (Promedio):} MAE = $16.25$ mmHg
        \item \textbf{Ridge:} MAE = $28.13$ mmHg
    \end{itemize}
\end{frame}

\begin{frame}{Resultados de Modelos Avanzados}
    \begin{block}{Sobreajuste (Overfitting)}
    Los algoritmos complejos no lograron generalizar con tan poca data, rindiendo por debajo del predictor Dummy.
    \end{block}
    
    \begin{itemize}
        \item \textbf{Random Forest:} MAE = $19.63$ mmHg
        \item \textbf{LightGBM:} MAE = $19.64$ mmHg
        \item \textbf{MLP (PyTorch):} MAE = $32.65$ mmHg
    \end{itemize}
\end{frame}

\begin{frame}{Siguientes Pasos}
    \begin{itemize}
        \item \textbf{Volumetría:} Procesar el conjunto completo de PulseDB (5,361 pacientes) y no solo la muestra exploratoria de 41 pacientes.
        \item \textbf{Ablation:} Continuar eliminando características morfológicas correlacionadas.
        \item \textbf{Tuning:} Habilitar la búsqueda Optuna en el servidor final.
    \end{itemize}
\end{frame}

\end{document}
